\section{The sharp Brun--Titchmarsh inequality}
\label{sec:sharp-bt}

We now supply the additional argument needed for the constant $2$.
Its essential analytic step is a finite inequality for points of a
circle with specified separation.  A rational grid replaces the
maximizing eigenvector that often appears in proofs of that inequality.
The subsequent estimates use finite divisor sums and explicit
logarithm bounds.  This follows the weighted large-sieve route to
Brun--Titchmarsh, with a local-spacing proof tailored to the Farey
fractions needed here. The classical weighted-sieve and Hilbert-inequality sources are \cite{MontgomeryVaughan1973,MontgomeryVaughan1974}; the finite-grid replacement and the explicit elementary completion are given in full below.

\begin{theorem}[Sharp Brun--Titchmarsh, arithmetical inputs]
\label{thm:sharp-bt}
For a rational $x>q$, an integer $q\geq1$, and $\gcd(a,q)=1$,
\begin{equation}
 \boxed{\displaystyle
 \pi(x;q,a)\leq\frac{2x}{\varphi(q)\log(x/q)}.}
 \label{eq:sharp-bt}
\end{equation}
\end{theorem}

Here $\pi(x;q,a)$ counts primes up to $\lfloor x\rfloor$.
All logarithms in the proof have the rational-approximation meaning
already specified.  Thus \eqref{eq:sharp-bt} means that, for every
precision parameter $r$, the corresponding rational lower
approximation $L_r(x/q)$ satisfies
$\varphi(q)\pi(x;q,a)L_r(x/q)\leq2x$.
The finite numerical side conditions below are checked using exact
integers and rational enclosures; their reproducible certificate is
included with the archive.  This numerical certificate checks the
side conditions, not the universal proof or an iPRA derivation.

\subsection{Three estimates from disjoint intervals}

Let $x_1,\ldots,x_R$ be rational points distinct modulo $1$, with
$R\geq2$.  Suppose positive rational numbers $\delta_r$ satisfy
\begin{equation}
 \delta_r\leq\min_{s\ne r}\|x_r-x_s\|,
 \qquad \delta_r\leq\frac18.
 \label{eq:local-spacing}
\end{equation}
Here $\|t\|$ is the distance to the nearest integer.  Write
\[
 K_{rs}=\csc\bigl(\pi(x_r-x_s)\bigr),\qquad
 C_{rs}=\cot\bigl(\pi(x_r-x_s)\bigr)
 \quad(r\ne s),
\]
and set both diagonal entries to zero.  Intervals $I_r$ of length
$\delta_r$, centered at $x_r$ on the circle, have disjoint interiors:
the center distance is at least each of $\delta_r,\delta_s$, hence at
least $(\delta_r+\delta_s)/2$.

The functions $\csc^2(\pi t)$ and $\csc^4(\pi t)$ are convex between
their poles.  Their products with translated copies of
$\csc^2(\pi t)$ are also convex there: the logarithmic second derivative
of each positive factor is positive.  Applying the midpoint integral
inequality on the intervals $I_r$ gives
\begin{align}
 \sum_{r\ne s}\delta_r K_{rs}^2
 &\leq\frac4{\pi^2\delta_s},
 \label{eq:spacing-two}\\
 \sum_{r\ne s}\delta_r K_{rs}^4
 &\leq\frac6{\pi^4\delta_s^3},
 \label{eq:spacing-four}\\
 \sum_{r\ne s,t}\delta_r K_{rs}^2K_{rt}^2
 &\leq\frac4{\pi^2}
 \left(\frac1{\delta_s}+\frac1{\delta_t}\right)K_{st}^2
 \qquad(s\ne t).
 \label{eq:spacing-mixed}
\end{align}
We verify the constants and the sign needed in the third estimate.

For \eqref{eq:spacing-two}, translate $x_s$ to zero and put
$b=\delta_s/2$.  The sum is bounded by
\[
 \int_b^{1-b}\csc^2(\pi t)\,dt
 =\frac2\pi\cot(\pi b)
 \leq\frac4{\pi^2\delta_s},
\]
using $\cot u\leq1/u$ for $0<u\leq\pi/2$.  For the fourth powers,
the same integral calculation gives
\[
 \frac2\pi\left(\cot(\pi b)+\frac13\cot^3(\pi b)\right)
 \leq\frac4{\pi^2\delta_s}
       +\frac{16}{3\pi^4\delta_s^3}.
\]
After multiplication by $\pi^4\delta_s^3$, its right side is at most
$16/3+4\pi^2\delta_s^2<143/24<6$, since
$\delta_s\leq1/8$ and $\pi^2<10$.  This proves
\eqref{eq:spacing-four}.

For \eqref{eq:spacing-mixed}, take $x_s=0$, $x_t=a\in(0,1)$,
$b=\delta_s/2$, and $c=\delta_t/2$.  The sum is at most the integral
of $\csc^2(\pi t)\csc^2(\pi(t-a))$ over
$[b,a-c]\cup[a+c,1-b]$.  The identity
\begin{align*}
 &\csc^2(\pi t)\csc^2(\pi(t-a))\\
 &\quad=\csc^2(\pi a)
 \left(\csc^2(\pi t)+\csc^2(\pi(t-a))
       +2\cot(\pi a)\bigl(\cot(\pi t)-\cot(\pi(t-a))\bigr)\right)
\end{align*}
separates the integral into three parts.  The first two have the
bounds in \eqref{eq:spacing-two}.  The integral of the difference of
cotangents is
\[
 \frac1\pi\log
 \frac{\sin(\pi(a-b))\sin(\pi(a-c))}
      {\sin(\pi(a+b))\sin(\pi(a+c))}.
\]
All four sine factors are positive.  If $a\leq1/2$, each numerator
factor is at most its matching denominator factor, whereas
$\cot(\pi a)\geq0$.  If $a\geq1/2$, both comparisons reverse.
In both cases the cotangent times this logarithm is nonpositive.
Dropping that term proves \eqref{eq:spacing-mixed}.  The comparison of
the paired sine factors is the identity
$\sin(\pi(a+h))-\sin(\pi(a-h))
=2\cos(\pi a)\sin(\pi h)$.

\subsection{A weighted cosecant inequality, with finite maximization}

\begin{lemma}[Local weighted cosecant inequality]
\label{lem:weighted-cosecant}
Under \eqref{eq:local-spacing}, for every finite complex vector $u$,
\begin{equation}
 \left|\sum_{r\ne s}u_r\overline{u_s}K_{rs}\right|
 \leq\frac32\sum_r\frac{|u_r|^2}{\delta_r}.
 \label{eq:weighted-cosecant}
\end{equation}
\end{lemma}

\begin{proof}
We first display the finite algebra behind the estimate, and then
give the rational search that supplies the approximate vector used
in that algebra.  Put $D=\operatorname{diag}(\delta_r)$ and
$H=-iD^{1/2}KD^{1/2}$.  The matrix $H$ is Hermitian.
Since $|\sin(\pi t)|\geq2\|t\|$ and
$\sqrt{\delta_r\delta_s}\leq\|x_r-x_s\|$, every entry of $H$ has
absolute value at most $1/2$.  In particular, the finite Cauchy
inequality gives $\|Hv\|\leq R\|v\|$.

Suppose temporarily that
\[
 \sum_r\frac{|u_r|^2}{\delta_r}=1,
 \qquad Ku=i\mu D^{-1}u+e,
 \qquad \sum_r\delta_r|e_r|^2\leq\epsilon^2,
\]
where $u=D^{1/2}v$, $\|v\|=1$, and
$\mu=v^*Hv\in\mathbb R$.  If $w=Hv-\mu v$, then $v^*w=0$ and
$e=iD^{-1/2}w$.  Hence
\begin{equation}
 \mu^2\leq\sum_r\delta_r|(Ku)_r|^2=T_1+T_2,
 \label{eq:cosecant-expansion}
\end{equation}
where
\[
 T_1=\sum_s|u_s|^2\sum_{r\ne s}\delta_r K_{rs}^2
 \leq\frac4{\pi^2}
\]
and
\[
 T_2=\sum_{\substack{r,s,t\text{ distinct}}}
       \delta_r u_s\overline{u_t}K_{rs}K_{rt}.
\]

The identity $K_{rs}K_{rt}=K_{st}(C_{rs}-C_{rt})$ gives
\begin{equation}
 T_2=U-V+\sum_{s\ne t}(\delta_s+\delta_t)
             u_s\overline{u_t}K_{st}C_{st},
 \label{eq:cosecant-cross}
\end{equation}
where
\[
 U=\sum_{\substack{r\ne s\\s\ne t}}
       \delta_r u_s\overline{u_t}K_{st}C_{rs},
 \qquad
 V=\sum_{\substack{r\ne t\\s\ne t}}
       \delta_r u_s\overline{u_t}K_{st}C_{rt}.
\]
The $r=t$ terms added to $U$ and the $r=s$ terms added to $V$
account exactly for the last sum in \eqref{eq:cosecant-cross}.
Substituting $Ku=i\mu D^{-1}u+e$ shows that the $\mu$ terms in
$U$ and $V$ are equal.  Their difference is
\begin{equation}
 U-V=2\sum_{r\ne s}\delta_r C_{rs}
                     \operatorname{Re}(u_s\overline{e_s}).
 \label{eq:cross-residual}
\end{equation}
For example, both $\mu$ contributions equal
$-i\mu\sum_{r\ne s}(\delta_r/\delta_s)|u_s|^2C_{rs}$;
this checks the cancellation including its signs.

Let $B\geq1/\min_r\delta_r$ be an integer.  Then
$|C_{rs}|\leq B/2$, $\delta_r\leq1$, and
\[
 \sum_s|u_se_s|=\sum_s|v_sw_s|\leq\epsilon.
\]
Thus $|U-V|\leq RB\epsilon$.  Since $|C_{st}|\leq|K_{st}|$,
\eqref{eq:cosecant-cross} gives $T_2\leq E+RB\epsilon$, where
\[
 E=\sum_{s\ne t}(\delta_s+\delta_t)|u_su_t|K_{st}^2.
\]
Rewrite this as
$E=2\sum_r|u_r|\sum_{s\ne r}\delta_s|u_s|K_{rs}^2$.
Weighted Cauchy, followed by
\eqref{eq:spacing-four} and \eqref{eq:spacing-mixed}, yields
\begin{align*}
 E^2
 &\leq4\sum_{s,t}\delta_s\delta_t|u_su_t|
           \sum_{r\ne s,t}\delta_rK_{rs}^2K_{rt}^2\\
 &\leq\frac{24}{\pi^4}+\frac{16}{\pi^2}E.
\end{align*}
The diagonal $s=t$ contributes at most $6/\pi^4$ before the
outer factor $4$; the other terms contribute at most $4E/\pi^2$.
For $X=\pi^2E\geq0$, therefore, $X^2\leq24+16X$.
It follows that $X<35/2$, because
$(35/2)^2-16(35/2)-24=9/4>0$ and the quadratic is increasing
beyond $8$.  Consequently, if $RB\epsilon\leq1/1000$,
\begin{equation}
 \mu^2\leq\frac{43}{2\pi^2}+\frac1{1000}
 <\frac{2150}{961}+\frac1{1000}
 <\left(\frac{1499}{1000}\right)^2.
 \label{eq:approximate-spectral-bound}
\end{equation}
We used $\pi>31/10$.  Each last comparison is rational arithmetic.

It remains to produce a sufficiently good approximate maximizer
without an existence theorem for an eigenvector.  Choose
\[
 \epsilon=\frac1{1000RB},\qquad
 \tau=\frac{\epsilon^2}{100R},\qquad
 h=\frac{\tau}{100R^2}.
\]
Approximate $H$ by a Hermitian matrix $\widehat H$ with Gaussian
rational entries, each within $h$ of the corresponding entry of
$H$.  Enumerate the Gaussian rational grid of mesh $h$ in the
coordinate box $[-2,2]^{2R}$, retaining vectors with squared norm
between $1/4$ and $4$.  This is a nonempty finite set.  Maximize
the rational quantity
$w^*\widehat H w/(w^*w)$ by a finite comparison loop, and normalize
the selected vector to obtain $v$.

For any unit vector, rounding its coordinates to this grid changes
its norm by at most $2Rh$, and normalization at most doubles that
distance.  Even allowing an extra factor two for first approximating
the input coordinates, the Rayleigh quotient changes by at most
$16R^2h$.  The matrix approximation changes a normalized Rayleigh
quotient by at most $Rh$ at either endpoint.  The chosen values of
$h$ therefore give
\begin{equation}
 v^*Hv\geq u^*Hu-\tau
 \quad\text{for every unit vector }u.
 \label{eq:grid-maximizer}
\end{equation}
For the rounded vector the squared norm remains in $[1/4,4]$,
so it is among the enumerated candidates.

Put $\mu=v^*Hv$ and $w=Hv-\mu v$.  Then $v^*w=0$ and
$\|w\|\leq2R$.  Substitute
\[
 u=\frac{v+w/(2R)}{\|v+w/(2R)\|}
\]
into \eqref{eq:grid-maximizer}.  Expanding its Rayleigh quotient,
using $\|H\|\leq R$ and $|\mu|\leq R$, gives
\[
 u^*Hu-\mu\geq\frac{\|w\|^2}{4R}.
\]
Indeed the numerator of the difference is at least
$\|w\|^2/(2R)$, and its denominator
$1+\|w\|^2/(4R^2)$ is at most $2$.
Thus $\|w\|^2\leq4R\tau$, so $\|w\|\leq\epsilon/5<\epsilon$.
This is the approximate vector required above.

Equation \eqref{eq:approximate-spectral-bound}, followed by
\eqref{eq:grid-maximizer}, bounds every unit Rayleigh quotient by
$1499/1000+\tau<3/2$.  Applying the same construction to $-H$
bounds its negative.  Rescaling the vector gives
\eqref{eq:weighted-cosecant}.  The entire search is finite, with
mesh and approximation precision specified by rational formulas.
\end{proof}

\subsection{From the cosecant bound to a weighted sieve}

For a positive integer $N$ and an arbitrary integer $M$, expanding
a finite square sum and using the geometric identity
\[
 \sum_{n=M+1}^{M+N}e^{2\pi in\theta}
 =\frac{e^{2\pi i(M+N+1/2)\theta}
             -e^{2\pi i(M+1/2)\theta}}
        {2i\sin(\pi\theta)}
\]
shows that
\[
 \sum_{n=M+1}^{M+N}\left|\sum_rv_re^{2\pi inx_r}\right|^2
 \leq\sum_r\left(N+\frac3{2\delta_r}\right)|v_r|^2.
\]
The diagonal contribution is $N\sum|v_r|^2$; the other terms are a
half-difference of two cosecant forms, each bounded by
Lemma~\ref{lem:weighted-cosecant}.  Dividing matrix columns by the
displayed positive square roots and applying finite Cauchy duality
gives
\begin{equation}
 \sum_r\frac{\left|\sum_{n=M+1}^{M+N}a_ne^{2\pi inx_r}\right|^2}
                  {N+3/(2\delta_r)}
 \leq\sum_{n=M+1}^{M+N}|a_n|^2.
 \label{eq:weighted-large-sieve}
\end{equation}
For completeness, if a finite matrix $A$ satisfies
$\|Av\|\leq\|v\|$, then
$\|A^*w\|^2=\langle AA^*w,w\rangle
\leq\|A^*w\|\|w\|$, which gives the dual inequality.
This step has no spectral or compactness hypothesis.

Take all distinct reduced fractions $a/d$ modulo one with
$1\leq d\leq z$, where $z\geq8$ is an integer, including $0/1$.
Set $\delta_{a/d}=1/(dz)$.  Their circular distances are at least
$1/(de)$, so \eqref{eq:local-spacing} holds.

Let $\mathcal A$ be a set of $Z$ integers in an interval of $N$
consecutive integers.  Suppose it omits one residue class modulo
each prime $p\leq z$ not dividing $q$.  For squarefree
$d\leq z$ coprime to $q$, CRT supplies $b_d$ such that
$\gcd(n-b_d,d)=1$ for every $n\in\mathcal A$.  If
$S(t)=\sum_{n\in\mathcal A}e^{2\pi int}$, the finite Ramanujan-sum
identity gives
\[
 \sum_{\substack{a\bmod d\\\gcd(a,d)=1}}
      e^{-2\pi iab_d/d}S(a/d)=\mu(d)Z.
\]
To verify the identity, expand the sum over reduced residues by
M\"obius inclusion-exclusion and sum finite geometric progressions;
when $\gcd(n-b_d,d)=1$, only the divisor $1$ remains.
Cauchy therefore yields
$\sum_{\gcd(a,d)=1}|S(a/d)|^2\geq Z^2/\varphi(d)$.
Substitute these lower bounds into \eqref{eq:weighted-large-sieve},
discard other nonnegative terms, and divide by $Z$ when $Z>0$.
The case $Z=0$ is immediate.  We obtain
\begin{equation}
 Z\leq
 \left(\sum_{\substack{d\leq z\\\gcd(d,q)=1}}
       \frac{\mu(d)^2}{\varphi(d)(N+3dz/2)}\right)^{-1}.
 \label{eq:weighted-arithmetic-sieve}
\end{equation}

For any positive decreasing weight $f$, a useful comparison is
\begin{equation}
 \sum_{d\leq z}\frac{\mu(d)^2f(d)}{\varphi(d)}
 \leq\frac q{\varphi(q)}
       \sum_{\substack{d\leq z\\\gcd(d,q)=1}}
                  \frac{\mu(d)^2f(d)}{\varphi(d)}.
 \label{eq:weighted-remove-modulus}
\end{equation}
Write each squarefree $d$ as $er$, with $e\mid\operatorname{rad}(q)$
and $\gcd(r,q)=1$.  Then $f(er)\leq f(r)$ and $r\leq z/e\leq z$.
Summing $1/\varphi(e)$ over these divisors gives $q/\varphi(q)$.
This proves \eqref{eq:weighted-remove-modulus} by a finite calculation.

\subsection{An explicit lower bound for the denominator}

Write $G(u)=\sum_{d\leq u}\mu(d)^2/\varphi(d)$, with real $u\geq1$.
We require
\begin{equation}
 G(u)>\log u+\frac{107}{100}\qquad(u\geq6).
 \label{eq:strong-G-bound}
\end{equation}
Here is a proof with a small finite certificate.

Define $\gamma$ by the monotone brackets
$H_n-\log(n+1)\leq\gamma\leq H_n-\log n$.
Their gap is at most $1/n$, giving an explicit Cauchy modulus.
For $y\geq1$, $m=\lfloor y\rfloor$, these brackets imply
\begin{equation}
 |H_m-\log y-\gamma|\leq\frac1y.
 \label{eq:harmonic-gamma-error}
\end{equation}
For the upper inequality, use
$\gamma\geq H_m-\log(m+1)$ and
$\log((m+1)/y)\leq(m+1-y)/y\leq1/y$.
For the lower one use
$\gamma\leq H_{m+1}-\log(m+1)$; the result is at least
$-1/(m+1)\geq-1/y$.

Let $H_6(y)=\sum_{n\leq y,\,\gcd(n,6)=1}1/n$.
Inclusion-exclusion over $1,2,3,6$ and
\eqref{eq:harmonic-gamma-error} give, for $y\geq6$,
\[
 H_6(y)=\frac13\log y+\frac13\gamma
       +\frac13\log2+\frac16\log3+E(y),
 \qquad |E(y)|\leq\frac4y.
\]
Decomposing squarefree integers according to their common divisor
with $6$, and using the radical bound of the preceding section,
gives, for $u\geq36$,
\begin{equation}
 G(u)\geq\log u+\gamma+\frac12\log2+\frac16\log3-\frac{30}{u}.
 \label{eq:G-large-bound}
\end{equation}
The coefficients follow from
$\sum_{m\mid6}1/\varphi(m)=3$ and
$\sum_{m\mid6}m/\varphi(m)=15/2$.

The certified rational inequalities
$\gamma>577/1000$, $\log2>692/1000$, and
$\log3>1098/1000$ show that
\eqref{eq:G-large-bound} implies \eqref{eq:strong-G-bound} for
$u\geq1000$.  For $6\leq u<1000$, let $n=\lfloor u\rfloor$.
The $994$ finite inequalities
\begin{equation}
 G(n)>\log(n+1)+\frac{107}{100},
 \qquad n=6,\ldots,999,
 \label{eq:G-finite-certificate}
\end{equation}
prove the result on each interval $[n,n+1)$.
They are verified using exact rational upper bounds for
$\log(n+1)$.  The smallest certified lower bound for
$G(n)-\log(n+1)$ exceeds $1.114081573$, at $n=9$.
The certificate for $\gamma>577/1000$ is
$H_{10000}-\operatorname{upperlog}(10001)>577/1000$;
its certified lower bound exceeds $0.577165669$.

The logarithm certificates use binary range reduction to $[1,2]$
and $18$ terms of
$2\sum_{j\geq0}w^{2j+1}/(2j+1)$, with its exact geometric-tail
bound.  The full rational verification program accompanies this
article.  Thus these are finite comparisons with proved error
bounds, rather than floating-point tests.

For $z\geq100$, define
\[
 J(z)=\sum_{d\leq z}\frac{\mu(d)^2}{\varphi(d)(1+d/z)}.
\]
Partial summation, using \eqref{eq:strong-G-bound} above $6$ and
retaining the terms $d\leq5$ exactly, gives, with $C=107/100$,
\begin{align*}
 J(z)\geq{}&
 \sum_{d\leq5}\frac{\mu(d)^2}{\varphi(d)}
              \left(\frac z{z+d}-\frac z{z+6}\right)\\
 &+\frac z{z+6}(\log6+C)
   +\int_6^z\frac z{u(z+u)}\,du.
\end{align*}
The integral equals
$\log z+\log(1+6/z)-\log12$.
Moreover,
$\sum_{d\leq5}\mu(d)^2(6-d)/\varphi(d)=43/4$,
$z/(z+d)\geq9/10$, and
$\log(1+6/z)\geq6/(z+6)$.
Consequently
\[
 J(z)-\log z\geq C-\log2+
             \frac{627/40-6(\log6+C)}{z+6}.
\]
The rational bounds $\log2<347/500$, $\log6<224/125$ and
$z+6\geq106$ make the right side at least
$38359/106000>361/1000$.  We have proved
\begin{equation}
 J(z)>\log z+\frac{361}{1000}\qquad(z\geq100).
 \label{eq:weighted-J-bound}
\end{equation}

We also use
\begin{equation}
 \pi(z)\leq z/3\qquad\text{for integers }z\geq100.
 \label{eq:weak-prime-count}
\end{equation}
For $z\geq165$, inclusion-exclusion with $2,3,5$ gives
$\pi(z)\leq4z/15+11\leq z/3$.
The $65$ integers $100\leq z<165$ are checked exactly in the
same finite certificate.

\subsection{The large range}

Let $t=x/q>16000$.  Choose the representative $a_0$ in
$\{1,\ldots,q\}$ and let $N$ be the number of integers
$a_0+qj\leq x$, $j\geq0$.  Then
$N-1\leq t<N+1$ and $N>15000$.
Take the integer
\[
 z=\max\{n\leq N:3n^2\leq2N\}
   =\left\lfloor\sqrt{2N/3}\right\rfloor\geq100.
\]
The indices of the primes exceeding $z$ omit one residue modulo
each prime $p\leq z$ not dividing $q$.  Apply
\eqref{eq:weighted-arithmetic-sieve} and
\eqref{eq:weighted-remove-modulus}.  Since
$3z^2/2\leq N$, we have
$N+3dz/2\leq N(1+d/z)$, and hence
\begin{equation}
 \pi(x;q,a)\leq\frac q{\varphi(q)}
                \left(\pi(z)+\frac N{J(z)}\right).
 \label{eq:large-range-sieve}
\end{equation}

The other root inequality, $(z+1)^2>2N/3$, gives
\[
 \log z\geq\tfrac12\log N-\tfrac12\log(3/2)-1/z.
\]
Using $\log(3/2)<203/500$ and \eqref{eq:weighted-J-bound},
\[
 J(z)>\tfrac12\log N+\tfrac{361}{1000}
                        -\tfrac{203}{1000}-\tfrac1{100}
       >\tfrac12\log N+\tfrac7{50}.
\]
Set $u=\log N$.  By \eqref{eq:weak-prime-count},
\[
 \pi(z)+N/J(z)\leq\frac{\sqrt N}{3}
                       +\frac{2N}{u+7/25}.
\]
For $N\geq15000$, we have $u\geq3$ and
\[
 \frac{2N}{u}-\frac{2N}{u+7/25}\geq\frac N{2u^2}.
\]
Also $u^2\leq\sqrt N$: if $w=N^{1/4}\geq10$, then
$\log w\leq w/4$.  This holds at $10$ because
$\log10<5/2$, and it persists thereafter by integrating
$1/w\leq1/4$.  Thus $u=4\log w\leq w$.
Finally $u\sqrt N\geq12$.  Combining these inequalities gives
\[
 \pi(z)+N/J(z)
 \leq\frac{2N}{u}-\frac{\sqrt N}{6}
 \leq\frac{2(N-1)}{\log N}.
\]
If $N-1\leq t\leq N$, then
$(N-1)/\log N\leq t/\log t$ by separately comparing numerator
and denominator.  If $t\geq N$, use the monotonicity of
$t/\log t$ for $t\geq e$.  Equation
\eqref{eq:large-range-sieve} now proves \eqref{eq:sharp-bt}
throughout $t>16000$.

\subsection{The small range}

Finite inclusion-exclusion gives, for any fixed sieving endpoint $z$,
\[
 \pi(x;q,a)\leq\frac{x}{\varphi(q)}
 \left(\prod_{p\leq z}(1-1/p)
       +\frac{2^{\pi(z)}+\pi(z)}{t}\right).
\]
For the main term remove primes dividing $q$ and use
$\prod_{p\mid q}(1-1/p)^{-1}=q/\varphi(q)$.
The error counts at most $2^{\pi(z)}$ divisor terms, and the
discarded small primes contribute at most $\pi(z)$.

Taking $z=2,3,5,13$, it is enough that one of the following four
functions is less than one:
\[
\begin{array}{ll}
 f_1(t)=\log t\left(\frac14+\frac3{2t}\right),&
 f_2(t)=\log t\left(\frac16+\frac3t\right),\\[3pt]
 f_3(t)=\log t\left(\frac2{15}+\frac{11}{2t}\right),&
 f_4(t)=\log t\left(\frac{96}{1001}+\frac{35}t\right).
\end{array}
\]
The respective ranges are
\[
 (1,20],\qquad[18,270],\qquad[42,1000],\qquad[700,20000].
\]
They overlap and cover $(1,20000]$.

For $f(t)=\log t(A+B/t)$, the sign of $f'$ is the sign of
$g(t)=At+B(1-\log t)$.
For $f_1$, the minimum of $g$ is attained at $6$ and equals
$3-(3/2)\log6>0$, so $f_1$ is increasing.
For each of the other intervals, its left endpoint is at least
$B/A$.  There $g$ is increasing, so $f$ has no interior local
maximum and is bounded by its endpoint values.
The finite certificate verifies
\[
 f_1(20)<1,\quad f_2(18),f_2(270)<1,\quad
 f_3(42),f_3(1000)<1,\quad f_4(700),f_4(20000)<1
\]
using rational upper logarithm bounds.  The smallest certified
endpoint margin is greater than $0.004724985$.
This proves \eqref{eq:sharp-bt} for $1<t\leq20000$, overlapping
the large range and completing the proof.

\subsection{Primitive recursive implementation and Lean obligations}

The proof uses no unbounded search for a prime.  Trial division,
totients, M\"obius values, CRT representatives, finite Fourier sums,
the integer square root, grid enumeration, and exact comparison of
rational Rayleigh quotients all have explicit primitive recursive
implementations.  The numerical certificate uses finite loops with
the stated bounds $10000$, $1000$, and $165$.

For the Farey set, the spacing data are rational:
$\delta_{a/d}=1/(dz)$ and $B=z^2$ suffice.  The number of points is
at most $z^2+1$.  Thus the matrix grid parameters in the proof are
explicit rational functions of $z$ and that point count.  Square
roots of positive rational numbers use bounded bisection.
Sine, cosine, and the exponential on bounded arguments use their
factorial series.  Their reciprocal denominators are apart from
zero by $|\sin(\pi(x_r-x_s))|\geq2/z^2$, which supplies a
primitive recursive precision bound for every division.

The integrations in the spacing estimates range over intervals at
distance at least $1/(2z^2)$ from their poles.  On these intervals,
the derivatives appearing in the midpoint and integration identities
have explicit bounds obtained by differentiating rational expressions
in sine and cosine.  Uniform Riemann-sum error bounds therefore
reduce the integral inequalities to rational finite sums at a
primitive recursive mesh.  The constants $\pi>31/10$ and
$\pi^2<10$ have finite alternating-series certificates using
$100$ terms of $\pi=4\arctan(1)$ and the next-term error bound.

A Lean development can follow the proof in the following order:
certified elementary functions; convex midpoint estimates for the
three explicit kernels; rational finite-grid maximization;
the residual cancellation \eqref{eq:cross-residual}; the weighted
cosecant and large-sieve inequalities; the Ramanujan-sum identity;
the decreasing-weight comparison; the exact side-condition
certificate; and the two final ranges.  The finite grid is essential
to the foundational audit: replacing it by an unrestricted spectral
theorem would lose the explicit constructive justification supplied
here.  No completed Lean or object-level iPRA proof is asserted by
this mathematical reconstruction.
